% ============================================================================
%  Fundamentação Teórica
%  Abertura da seção: o que os três métodos têm em comum e o que os diferencia.
%  Em seguida, cada método é desenvolvido na ordem:
%     (1) descrição analítica geral  (2) aplicação aos parâmetros fixos
%     (3) reconstrução               (4) gancho para o método seguinte
% ============================================================================

\section{Fundamentação Teórica}

\begin{frame}{Amostragem}{Estrutura Comum aos Três Métodos}
    \small
    Amostrar é registrar o sinal contínuo apenas nos instantes $t = nT_s$. Os três
    métodos partem do mesmo trem de impulsos e diferem somente no \textbf{pulso} que
    carrega cada amostra:
    \begin{equation*}
        x_{am}(t) = \underbrace{\left[\, x_f(t) \cdot \sum_{n=-\infty}^{\infty} \delta(t - nT_s) \right]}_{\text{comum aos três métodos}}
        \;\longrightarrow\; \text{pulso de cada método}
    \end{equation*}

    \begin{center}
    \begin{tabular}{lll}
        \hline
        \textbf{Método} & \textbf{Amostra transportada por} & \textbf{Consequência} \\
        \hline
        Ideal & impulso $\delta(t - nT_s)$ & recuperação exata \\
        Natural & pulso que \textit{acompanha} $x_f(t)$ & espectro ponderado por $\mathrm{sinc}$ \\
        Topo-plano & pulso de topo constante $h(t)$ & $\mathrm{sinc}$ + atraso $\tau/2$ \\
        \hline
    \end{tabular}
    \end{center}

    \vspace{0.5em}
    Começamos pelo \textbf{ideal}: é o caso sem pulso, e portanto a referência com que
    os outros dois serão comparados.
\end{frame}
